The final linguistic correction is the most time-consuming step in the completion
of a paper and will probably require several passes.
Nevertheless, this step should not be approached half-heartedly.
An accumulation of typing errors rightly raises the question of whether the
author has exercised the necessary thoroughness in terms of content.
There are several ways to check spelling, grammar and expression. As a rule
a combination of all variants leads to the goal.


\paragraph{Automatic spell check}

The automatic spell checker is the simplest method of spelling correction.
It is the best way to detect typing errors.
Many modern \LaTeX editors offer corresponding functions. A normal text editor
may be confused by \LaTeX commands. Some editors not only recognise typing
errors, but also grammatically questionable constructions or highlight
repetitions and the use of redundant filler words.

If you are writing your work with the command line under UNIX-like systems
the software Aspell\footnote{\url{http://aspell.net/}} offers the possibility of
software offers the possibility to check \TeX files. The command line call for
this is:

\begin{lstlisting}[escapechar=\!]
 aspell -t -c !\lara{filename}!
\end{lstlisting}

However, such tools do not detect bad expression or sentence structure.


\paragraph{Manual correction of a proof}

Make an early proof of chapters that have already been completed (except for
minor details). Sometimes \LaTeX commands in the manuscript prevent a direct
view of the final text. A look on the finished result helps to correct the
text for sentence structure and grammar.

It is easier to spot typing and spelling errors if you read each page
backwards. In this way, you avoid skipping words while reading, or read
differently from the way they appear on paper.

Cross out errors you recognise immediately. Do not cross out the error marks
in a different colour, but only after the error has been corrected in the \TeX
file. Do not throw away a page of the test print until all errors have been
crossed out.


\paragraph{Correction by uninvolved acquaintances}

Ask your acquaintances if someone would be willing to proofread your work. 
It can even be helpful if the proofreader is not a specialist. He or she will
then inevitably focus on the spelling and less on the content.

In order to be able to add comments to the manuscript more quickly, it is
advisable to number the lines in the proof copy. It is much easier to correct
word~five in line~256 than to correct word~five in the fourth~line
of paragraph~3 in section~2.3. \LaTeX\ inserts line numbers fully
automatically, by commenting out the lines given in Listing~\ref{lst:final:lineno}
in the template.

\lstinputlisting[language={[LaTeX]{TeX}},style=block,float,caption={turn on line numbering},label={lst:final:lineno}]{code/final_lineno.tex}

Do not forget to reinsert the comment characters before handing in the work.


\paragraph{Supervisor's feedback}

If this is a final paper, a good supervisor will request a preliminary version.
Depending on the workload and the job situation, the supervision in practice
may vary in intensity. With a little tact, however, you can persuade a
supervisor to read your work. However, please bear in mind that the mills of a
university grind slowly. Your supervisor certainly still has applications to
write in parallel, and conferences to attend at the same time. The supervisor
is not your secretary. He or she will certainly not tell you about typing and
questions about technical terms and the coherence of the content are clearly in
the foreground here.
